% ====================================================================
% SM2B_ensemble_equiv.tex  —  v1.0.22 창 SM2 / 신규 심화 후보 B (차선)
%   초안 전용 — 마스터 소스 미수정. 삽입처 = _sections/ch1_sec02b_part0.tex
%   §2.5(sec:sm-mc), eq:sm-mc-fluc verifybox 직후·소절 맺음 문단 직전
%   (현행 380행 "\end{verifybox}" 뒤, 382행 "이로써 사다리에는…" 앞).
%   목적 = §2.3 이 "앙상블 동등성"이라 명명만 하고 §2.5 가 대정준→정준 Legendre
%          반전이라 부른 것을, 방금 계산한 var(N) 로 물리(상대 요동 1/√M 소멸)로
%          닫아 두 앙상블이 같은 열역학을 주는 근거·반전 유일근의 거시 정당성을 준다.
%   ★절대 제약: 새 피팅 파라미터 0 · 기존 수식/라벨/기호 무변경 ·
%          제거 용이 독립 블록(무라벨 식) · 표준 정리(앙상블 동등성) 무인용 관용.
%   참조 라벨(전부 기존·같은 빌드): eq:sm-mucount · eq:sm-mc-fluc · eq:sm-mc-occ ·
%          eq:sm-mc-factor · eq:sm-thresh · sec:sm-lattice · sec:sm-mc · sec:hys.
%   서지 = hill1960 · mcquarrie1976 (V1·§2.5 에 이미 인용 — 신규 아님).
%   자산 앵커(제안): [V22-SM2-B]
% ====================================================================

\begin{bgbox}
\textbf{앙상블 동등성 --- 상대 요동의 $1/\sqrt{M}$ 소멸.} \S\ref{sec:sm-lattice} 는 대정준 경로(자리 하나의
확률)와 셈법 경로(배치 수의 로그)가 같은 $\mu(\theta)$ 에 닿는 것을 \emph{앙상블 동등성}이라 불렀고
(식~\eqref{eq:sm-mucount}), 이 소절은 고정 함량 반전을 대정준$\to$정준 Legendre 반전으로 닫았다. 두 앙상블이
같은 열역학을 주는 근거가 방금 계산한 요동에 있다. 식~\eqref{eq:sm-mc-fluc} 로 입자수 분산은 자리 수에
비례하고($\mathrm{var}(N)=\sum_jM_j\theta_j(1-\theta_j)=O(M)$, $M\equiv\sum_jM_j$) 평균도 자리 수에
비례하므로($\langle N\rangle=\sum_jM_j\theta_j=O(M)$, 식~\eqref{eq:sm-mc-occ}), \emph{상대} 요동은
\[
\frac{\sqrt{\mathrm{var}(N)}}{\langle N\rangle}\ \sim\ \frac{1}{\sqrt{M}}
\ \xrightarrow[\ M\to\infty\ ]{}\ 0
\]
로 거시 자리 수 $M$(마이크론 흑연 입자에서 $M\gg10^{8}$)에서 소멸한다. 곧 $\mu$ 를 고정하면 --- 대정준에서
$\langle N\rangle$ 이 요동하더라도 --- 함량이 사실상 한 값에 못박혀, $N$ 을 고정한 정준 기술과 같은 강도적
열역학을 준다. 대정준($\mu$ 독립)으로 세든 정준($\bar x$ 고정)으로 세든 결과가 같다는 앙상블 동등성이 이
소멸의 따름이며 \cite{hill1960,mcquarrie1976}, 이 소절의 반전이 대정준 근 $U_\mathrm{oc}$ 를 정준 관측(고정
$\bar x$)으로 읽어도 되는 근거다 --- $\sqrt{\mathrm{var}(N)}/\langle N\rangle$ 이 곧 $U_\mathrm{oc}$ 를
흐리는 함량 폭인데, 그 폭이 거시 $M$ 에서 소멸하기 때문이다. \emph{한정} --- 이 소멸은 자리
독립(식~\eqref{eq:sm-mc-factor})의 인수분해 전제 위에서다; 평균장 이중웰($\Omega_j>2RT$,
식~\eqref{eq:sm-thresh})의 상분리 창은 강상관 영역이라 요동이 거시적으로 자라 상대 요동이 $1/\sqrt{M}$ 로 죽지
않으며, 그 창의 분기 구성은 \S\ref{sec:hys} 소관이다(본론과 Part T 가 반전에 실제로 쓰는 전이별 강단조
서식은 이 소멸 안).
\end{bgbox}


% ─────────────────────────────────────────────────────────────────────
% 설계 주석
%   · 판형 = bgbox(배경 = "왜 두 앙상블이 같은 답인가"). 무라벨 식 1개 →
%     §2.5 기존 식번호(eq:sm-mc-*) 재배열 0 · 통삭제 시 소절 원상 복귀.
%   · §2.3 "앙상블 동등성"(명명)·§2.5 "Legendre 반전"(명명)에 메커니즘을 부여 →
%     중복 아님(둘 다 결론만 있었고 1/√M 논거는 전문에 전무 — grep 확인).
%   · var(N) 을 바로 위 eq:sm-mc-fluc 에서 물려받아 자족 — 새 계산 없음.
%   · 가드(Ω>2RT 강상관 예외)는 §2.5 verifybox 이중웰 가드·§4 소관과 정합.
%   · 표준 캡스톤(Hill Ch.그랜드캐노니컬 요동 / McQuarrie 앙상블 동등성)이라
%     hill1960·mcquarrie1976 앵커만 — 신규 서지 0.
% ─────────────────────────────────────────────────────────────────────
